The lithium battery industry invests heavily in ultra-high purity precursors, assuming that all trace impurities degrade performance. We challenge this assumption through a systematic computational study combining physics-based battery simulations (PyBaMM) with analysis of published experimental datasets. We model the effects of four common trace impurities---magnesium, water, iron, and oxygen-containing species---on Coulombic Efficiency (CE) and capacity retention across concentration ranges from 50 to 1{,}000~ppm. Our simulations reveal that magnesium at ${\sim}1{,}000$~ppm significantly improves both CE ($+0.0018\%$, $p=0.018$) and capacity retention ($+0.10\%$) by promoting a denser, lower-resistivity solid-electrolyte interphase (SEI). Oxygen-containing species show similar benefits at 200~ppm by promoting Li$_2$O-rich SEI, consistent with recent experimental findings. In contrast, iron is universally detrimental ($-0.33\%$ capacity retention at 1{,}000~ppm), and water exhibits a threshold effect---neutral below ${\sim}260$~ppm but degrading above 500~ppm. SEI resistivity emerges as the dominant predictor of impurity benefit (Pearson $r=-0.891$, $p<0.0001$). These results provide a simple design rule: impurities that reduce SEI resistivity are beneficial. Our findings suggest that controlled introduction of specific trace elements could simultaneously reduce manufacturing costs and improve battery performance, opening a new avenue for SEI engineering through deliberate ``impurity doping.''
